\documentclass[exercice]{thedocument}%
\usepackage{texmaths-tikz}%
\startdatakeys
\source{originale}
\BO{2026}
\cycle{4}
\competencesDuSocle{CA3,RA1}
\objectifsApprentissage{C4A2ea,C4A3ca}
\motsCles{aire du triangle, expression littérale, équation}
\enddatakeys
\begin{document}%
Soit ABC un triangle rectangle et iscoèle en A dont l'aire vaut 18 cm$^2$.\par
\medskip\noindent
\begin{tikzpicture}[scale=2]
    \coordinate[label=below left:A](A)at(0,0);
    \coordinate[label=above left:B](B)at(0,1);
    \coordinate[label=below right:C](C)at(1,0);
    \draw(A)--(B)--(C)--cycle;
    \tkzMarkSegments[mark=s||](A,B A,C)
    \tkzMarkRightAngle[size=0.125](C,A,B)
\end{tikzpicture}\par\medskip\noindent
On rappelle la formule donnant l'aire d'un triangle~:
$$\text{Aire}=\text{base}\times\text{hauteur}\div2$$
\begin{enumerate}
    \item En notant $x$ la longueur du segment [AB], exprimer l'aire du triangle ABC
    en fonction de $x$.
    \item Déterminer la longueur du segment [AB].
\end{enumerate}
\end{document}
